\documentclass[a4paper]{article}

% Basic packages
\usepackage[fontset=fandol]{ctex}
\usepackage{amsmath, amssymb}
\usepackage{graphicx}
\usepackage[margin=2.5cm]{geometry}
\usepackage[most]{tcolorbox}
\usepackage{etoolbox}
\usepackage{listings}
\usepackage{hyperref}
\usepackage{booktabs}
\usepackage{subcaption}
\usepackage{float}
\usepackage{tikz}
\usetikzlibrary{positioning, arrows.meta, calc}
\IfFileExists{pgfplots.sty}{
    \usepackage{pgfplots}
    \pgfplotsset{compat=1.18}
}{}

% Highlight boxes
\newtcolorbox{knowledgebox}[1]{
    enhanced, colback=blue!5!white, colframe=blue!75!black, colbacktitle=blue!75!black,
    coltitle=white, fonttitle=\bfseries, title=#1, attach boxed title to top left={yshift=-2mm, xshift=2mm},
    boxrule=1pt, sharp corners
}

\newtcolorbox{importantbox}[1]{
    enhanced, colback=yellow!10!white, colframe=yellow!80!black, colbacktitle=yellow!80!black,
    coltitle=black, fonttitle=\bfseries, title=#1, sharp corners
}

\newtcolorbox{warningbox}[1]{
    enhanced, colback=red!5!white, colframe=red!75!black, colbacktitle=red!75!black,
    coltitle=white, fonttitle=\bfseries, title=#1, sharp corners
}

% Listings
\lstset{
    language=Python,
    basicstyle=\ttfamily\small,
    keywordstyle=\color{blue},
    stringstyle=\color{red!60!black},
    commentstyle=\color{green!60!black},
    breaklines=true,
    frame=single,
    numbers=left,
    numberstyle=\tiny\color{gray},
    captionpos=b,
    extendedchars=false
}

% Front-page metadata
\newcommand{\notetitle}{CS224N Lecture 3: Backpropagation and Neural Networks}
\newcommand{\noteauthors}{整理：AI Course Notes Contributors；授课：Chris Manning}
\newcommand{\notedate}{2024年4月9日}
\newcommand{\videochannel}{Stanford Online}
\newcommand{\videopublishdate}{2024-04-09}
\newcommand{\videoduration}{1:13:23}
\newcommand{\videourl}{https://www.youtube.com/watch?v=HnliVHU2g9U}
\newcommand{\videocoverpath}{}
\newcommand{\repourl}{https://github.com/hqhq1025/ai-course-notes}


% Footer with repo info
\usepackage{fancyhdr}
\pagestyle{fancy}
\fancyhf{}
\fancyfoot[C]{\thepage}
\fancyfoot[R]{\footnotesize\color{gray}\href{\repourl}{AI Course Notes}}
\renewcommand{\headrulewidth}{0pt}
\renewcommand{\footrulewidth}{0pt}

\begin{document}
\setlength{\emergencystretch}{2em}

\begin{titlepage}
\centering
{\Large 课程笔记\par}
\vspace{1.2cm}
{\huge\bfseries \notetitle\par}
\vspace{0.8cm}
{\large \noteauthors\par}
\vspace{0.3cm}
{\large \notedate\par}
\vspace{1.2cm}

\vspace{1.2cm}
\vfill
\begin{tcolorbox}[width=0.9\textwidth, colback=blue!3!white, colframe=blue!55!black, sharp corners]
\textbf{课程版本：} 2024 年中文讲义整理版\par
\textbf{笔记来源：} AI Course Notes Contributors\par
\textbf{本版处理：} 移除课件截图和对应图注，不包含 Stanford 原始课件、图片或视频。\par
\textbf{授权：} 笔记依 CC BY-NC-SA 4.0 分享；完整许可与改动记录见下一页。
\end{tcolorbox}
\end{titlepage}

\begin{center}
\fbox{\begin{minipage}{0.91\textwidth}
\small
\textbf{来源与许可}\\
本中文笔记由 AI Course Notes Contributors 整理，课程版本为 2024；源稿：\url{https://github.com/hqhq1025/ai-course-notes/tree/717e2df65c3d2eee593c171156a30d4f8f6812b9/cs224n/lecture03}。\\
笔记内容依 CC BY-NC-SA 4.0 分享：\url{https://creativecommons.org/licenses/by-nc-sa/4.0/}。\\
本副本的改动：移除 Stanford 原始课件截图与依赖这些截图的图注，移除课程视频封面/链接，并清理 TeX 源稿中的行尾空白后重新编译为可独立阅读的 PDF；同时加入本来源与许可说明。完整许可文本随本文件夹提供。Stanford 课件、课件图片和视频不包含在该授权内；请从对应年份 Stanford 课程页访问原始课件。
\end{minipage}}
\end{center}
\clearpage

\tableofcontents
\newpage

%% ========== Part 1: 引言与神经网络回顾 ==========

\section{引言：课程概览}

本节课是 Stanford CS224N（Natural Language Processing with Deep Learning）的第三讲，由 Chris Manning 教授讲授。课程的核心目标是让学生理解神经网络训练的数学原理——如何\textbf{手动推导梯度}（矩阵微积分）以及如何\textbf{用算法自动计算梯度}（反向传播算法）。


% Raster figure omitted from this study copy.



课程分为三个部分：


% Raster figure omitted from this study copy.



\begin{importantbox}{本讲核心学习目标}
\begin{enumerate}
    \item 理解神经网络中梯度的数学推导——矩阵微积分（Matrix Calculus）
    \item 掌握反向传播算法（Backpropagation）——链式法则的高效实现
    \item 建立``前向传播计算函数值、反向传播计算梯度''的直觉
\end{enumerate}
\end{importantbox}

\subsection{本章小结}

本讲聚焦于神经网络训练的数学基础。Manning 教授强调，虽然现代深度学习框架（如 PyTorch）已经自动化了梯度计算，但理解底层原理对于调试、设计新模型和理解训练行为至关重要。

%% ========== Part 2: 神经网络基础回顾 ==========

\section{神经网络基础回顾}

\subsection{从逻辑回归到神经网络}

上一讲介绍了神经网络的基本概念：每个神经元本质上是一个小型的逻辑回归单元。与传统机器学习中单个逻辑回归不同的是，神经网络将这些单元\textbf{级联}起来，形成多层结构。


% Raster figure omitted from this study copy.



\begin{importantbox}{神经网络的核心优势——表示学习}
神经网络最强大的特性在于\textbf{中间层的自组织}（self-organization of intermediate representations）。网络不需要人工设计特征，而是自动学习``什么样的中间表示对最终任务有用''。这正是神经网络在大多数场景下优于传统机器学习的根本原因。
\end{importantbox}

\subsection{层的矩阵表示}

当神经网络采用规则的层状结构时，每一层的计算可以用矩阵运算表示：


% Raster figure omitted from this study copy.



具体来说，一层神经网络的计算包含两步：

\textbf{第一步：线性变换}
$$z = Wx + b$$

其中：
\begin{itemize}
    \item $x \in \mathbb{R}^m$：输入向量
    \item $W \in \mathbb{R}^{n \times m}$：权重矩阵
    \item $b \in \mathbb{R}^n$：偏置向量
    \item $z \in \mathbb{R}^n$：线性变换输出
\end{itemize}

\textbf{第二步：非线性激活}
$$a = f(z)$$

激活函数 $f$ 逐元素应用：$f([z_1, z_2, z_3]) = [f(z_1), f(z_2), f(z_3)]$。

\subsection{NER 示例：命名实体识别}

Manning 教授使用一个具体的 NER（命名实体识别）任务来贯穿整个讲解：判断一个上下文窗口中心的词是否为地点名称。


% Raster figure omitted from this study copy.



该网络的完整计算过程：
\begin{enumerate}
    \item 将 5 个词的词向量拼接为输入 $x \in \mathbb{R}^{5d}$
    \item 经过隐藏层：$h = f(Wx + b)$
    \item 点积得到分数：$s = u^T h$
    \item 通过 sigmoid 得到概率：$J_t(\theta) = \sigma(s) = \frac{1}{1 + e^{-s}}$
\end{enumerate}

\begin{knowledgebox}{为什么用词嵌入而不是 one-hot？}
One-hot 向量维度极高（词表大小）且稀疏，无法表达词与词之间的语义关系。词嵌入（word embeddings）将每个词映射为一个低维稠密向量，使得语义相近的词在向量空间中距离也近。在深度学习中，$\theta$ 不仅包含权重矩阵等传统参数，\textbf{词向量本身也是需要学习的参数}。
\end{knowledgebox}

\subsection{本章小结}

神经网络通过将多个``小型逻辑回归''级联，自动学习中间表示。每一层的计算可以简洁地表示为矩阵乘法加偏置再加非线性激活：$a = f(Wx + b)$。接下来的问题是：如何训练这个网络？答案是梯度下降，而这需要我们计算梯度。

%% ========== Part 3: 激活函数 ==========

\section{激活函数：为什么需要非线性}

\subsection{从阈值函数到可微激活函数}

神经元模型的历史始于 McCulloch \& Pitts（1943）的阈值单元：


% Raster figure omitted from this study copy.



阈值函数的问题在于它是分段常数函数——斜率处处为零（除了不连续点），因此\textbf{无法提供梯度信息}来指导学习。梯度学习的核心思想是``找到哪里更陡峭，朝那个方向走''——如同滑雪，你需要坡度来决定方向。

\subsection{常见激活函数一览}


% Raster figure omitted from this study copy.



\textbf{Logistic（Sigmoid）函数}：
$$\sigma(z) = \frac{1}{1 + \exp(-z)}$$
输出范围 $(0, 1)$，常用于输出层（映射到概率）。缺点：输出始终为正，且在两端梯度趋近于零（梯度消失）。

\textbf{Tanh 函数}：
$$\tanh(z) = \frac{e^z - e^{-z}}{e^z + e^{-z}}$$
输出范围 $(-1, 1)$，是 sigmoid 的缩放平移版本：$\tanh(z) = 2\sigma(2z) - 1$。

\textbf{ReLU（Rectified Linear Unit）}：
$$\text{ReLU}(z) = \max(z, 0)$$
正半轴斜率恒为 1，计算极快。虽然负半轴``死区''（梯度为零）看似矛盾，但实践中效果极好——它提供了一种自然的稀疏激活和特化机制。

\begin{warningbox}{ReLU 的``死区''问题}
当 ReLU 神经元的输入持续为负时，梯度为零，该神经元``永久死亡''，无法再被激活。这就是所谓的 \textbf{dying ReLU} 问题。虽然在正常训练中这通常不是致命问题（网络中总有一部分神经元是活跃的），但在学习率过大或初始化不当时可能导致大量神经元死亡。
\end{warningbox}

\textbf{Leaky ReLU / Parametric ReLU}：在负半轴给予一个小斜率（如 0.01），避免完全死亡。Parametric ReLU 将负半轴斜率作为可学习参数。

\textbf{Swish 和 GELU}：近年来在 Transformer 模型中广泛使用。
$$\text{Swish}(x) = x \cdot \sigma(x), \qquad \text{GELU}(x) = x \cdot P(X \le x), \; X \sim \mathcal{N}(0,1)$$

两者形状类似：正半轴近似 $y = x$，负半轴有一小段``弯曲''区域提供梯度。

\subsection{为什么非线性不可或缺}


% Raster figure omitted from this study copy.



\begin{importantbox}{线性层的叠加 = 单层线性变换}
如果只做矩阵乘法而没有非线性：
$$h_2 = W_2(W_1 x) = (W_2 W_1) x = W'x$$
多层网络在表示能力上与单层完全等价！非线性激活函数是使神经网络成为\textbf{万能函数逼近器}（universal function approximator）的关键。
\end{importantbox}

\begin{knowledgebox}{线性网络在学习理论中的价值}
Manning 教授指出了一个有趣的现象：虽然线性神经网络在\textbf{表示能力}上没有优势（等价于单层线性变换），但在\textbf{学习动力学}上却有独特性质。理论界有相当多论文研究线性网络的学习行为，因为多层结构即使没有非线性也会影响优化路径。
\end{knowledgebox}

\subsection{本章小结}

激活函数是神经网络能够学习复杂非线性映射的关键。从历史上的阈值函数到 sigmoid/tanh，再到 ReLU 及其变体，直到 Transformer 时代的 Swish/GELU，激活函数的演进反映了对``梯度流通性''和``计算效率''的不断追求。核心原则是：保留足够的梯度信号使学习成为可能。

%% ========== Part 4: 随机梯度下降回顾 ==========

\section{梯度下降与梯度计算}

\subsection{随机梯度下降（SGD）回顾}

神经网络的训练依赖于\textbf{随机梯度下降}（Stochastic Gradient Descent）：


% Raster figure omitted from this study copy.



$$\theta^{\text{new}} = \theta^{\text{old}} - \alpha \nabla_\theta J(\theta)$$

其中：
\begin{itemize}
    \item $\theta$：模型的所有参数（包括权重矩阵 $W$、偏置 $b$、词向量等）
    \item $\alpha$：学习率（step size）
    \item $\nabla_\theta J(\theta)$：损失函数关于参数的梯度
\end{itemize}

核心问题：\textbf{如何计算 $\nabla_\theta J(\theta)$？} 有两种方法：
\begin{enumerate}
    \item \textbf{手动推导}（By hand）——矩阵微积分
    \item \textbf{算法自动计算}（Algorithmically）——反向传播算法
\end{enumerate}

本讲两种方法都会讲解。

\subsection{本章小结}

SGD 是训练神经网络的核心算法，其关键步骤是计算损失函数关于所有参数的梯度。在深度学习中，参数 $\theta$ 不仅包括传统的权重和偏置，还包括词嵌入等数据表示参数。

%% ========== Part 5: 矩阵微积分 ==========

\section{矩阵微积分：手动推导梯度}

\subsection{从单变量到多变量}

Manning 教授给出了一个关键的思维口诀：

\begin{importantbox}{矩阵微积分的核心口诀}
\textbf{``Multivariable calculus is just like single-variable calculus if you use matrices.''}（多变量微积分就是单变量微积分——只不过用矩阵代替标量。）
\end{importantbox}


% Raster figure omitted from this study copy.



回顾单变量微积分的基本例子：$f(x) = x^3$，导数 $f'(x) = 3x^2$。导数衡量的是\textbf{斜率}——输入的微小变化如何被放大到输出上。例如在 $x = 4$ 处，$f'(4) = 48$，意味着 $x$ 从 4 变到 4.01 时，$f(x)$ 大约增加 $48 \times 0.01 = 0.48$（即从 64 变为约 64.48）。

\subsection{梯度与 Jacobian 矩阵}

当函数有 $n$ 个输入和 1 个输出时（$f: \mathbb{R}^n \to \mathbb{R}$），梯度是一个向量：

$$\nabla f = \begin{pmatrix} \frac{\partial f}{\partial x_1} \\ \frac{\partial f}{\partial x_2} \\ \vdots \\ \frac{\partial f}{\partial x_n} \end{pmatrix}$$

当函数有 $n$ 个输入和 $m$ 个输出时（$f: \mathbb{R}^n \to \mathbb{R}^m$），偏导数构成一个 $m \times n$ 的矩阵，称为 \textbf{Jacobian 矩阵}：

$$\frac{\partial \mathbf{f}}{\partial \mathbf{x}} = \begin{pmatrix} \frac{\partial f_1}{\partial x_1} & \cdots & \frac{\partial f_1}{\partial x_n} \\ \vdots & \ddots & \vdots \\ \frac{\partial f_m}{\partial x_1} & \cdots & \frac{\partial f_m}{\partial x_n} \end{pmatrix}$$


% Raster figure omitted from this study copy.



\begin{knowledgebox}{Jacobian 在神经网络中的角色}
神经网络的每一层都是一个多输入多输出的函数。例如一个隐藏层 $h = f(Wx + b)$ 将 $m$ 维输入映射到 $n$ 维输出。该层的 Jacobian 就是一个 $n \times m$ 的矩阵，描述了每个输出维度对每个输入维度的敏感度。
\end{knowledgebox}

\subsection{链式法则的矩阵形式}

对于复合函数，链式法则变为 Jacobian 矩阵的\textbf{乘法}。如果 $h = f(g(x))$，则：

$$\frac{\partial h}{\partial x} = \frac{\partial h}{\partial g} \cdot \frac{\partial g}{\partial x}$$

这与单变量的链式法则 $\frac{dz}{dx} = \frac{dz}{dy} \cdot \frac{dy}{dx}$ 形式完全一致——只是标量乘法变成了矩阵乘法。


% Raster figure omitted from this study copy.



\subsection{逐元素激活函数的 Jacobian}

对于逐元素激活函数 $h = f(z)$（其中 $h_i = f(z_i)$），其 Jacobian 是一个\textbf{对角矩阵}：


% Raster figure omitted from this study copy.



$$\left(\frac{\partial \mathbf{h}}{\partial \mathbf{z}}\right)_{ij} = \begin{cases} f'(z_i) & \text{if } i = j \\ 0 & \text{otherwise} \end{cases}$$

即：
$$\frac{\partial \mathbf{h}}{\partial \mathbf{z}} = \text{diag}(f'(\mathbf{z}))$$

\begin{importantbox}{对角结构的直觉}
为什么是对角矩阵？因为激活函数是\textbf{逐元素}应用的：$h_i$ 只依赖于 $z_i$，不依赖于 $z_j$（$j \neq i$）。因此非对角元素（$i \neq j$）的偏导数为零。
\end{importantbox}

\subsection{线性层的 Jacobian}

对于线性变换 $z = Wx + b$：

$$\frac{\partial z}{\partial x} = W, \qquad \frac{\partial z}{\partial b} = I$$

类比单变量：$y = wx + b$ 的导数关于 $x$ 是 $w$，关于 $b$ 是 $1$。矩阵版本中，$w$ 变为 $W$，$1$ 变为单位矩阵 $I$。

对于点积 $s = u^T h$：

$$\frac{\partial s}{\partial h} = u^T$$


% Raster figure omitted from this study copy.



\subsection{NER 网络的完整梯度推导}

现在将所有 Jacobian 组合起来，对 NER 示例网络推导梯度。网络结构为：

$$s = u^T h, \quad h = f(z), \quad z = Wx + b$$


% Raster figure omitted from this study copy.



\textbf{计算 $\frac{\partial s}{\partial b}$}：

$$\frac{\partial s}{\partial b} = \frac{\partial s}{\partial h} \cdot \frac{\partial h}{\partial z} \cdot \frac{\partial z}{\partial b} = u^T \cdot \text{diag}(f'(z)) \cdot I$$

化简后：
$$\frac{\partial s}{\partial b} = u^T \circ f'(z)$$

其中 $\circ$ 表示 \textbf{Hadamard 乘积}（逐元素乘法）。

\begin{knowledgebox}{Hadamard 乘积}
Hadamard 乘积（记作 $\circ$ 或 $\odot$）是将两个相同维度的向量/矩阵\textbf{逐元素相乘}。不同于点积（结果是标量），Hadamard 乘积的结果维度与输入相同。在神经网络的梯度计算中，它出现在对角 Jacobian 乘以向量时——等价于两个向量的逐元素乘法。
\end{knowledgebox}

\textbf{上游梯度 $\delta$ 的概念}：

注意到计算 $\frac{\partial s}{\partial b}$ 和 $\frac{\partial s}{\partial W}$ 时，前两项是共享的：

$$\delta = \frac{\partial s}{\partial h} \cdot \frac{\partial h}{\partial z} = u^T \circ f'(z)$$

$\delta$ 称为\textbf{上游梯度}（upstream gradient）或\textbf{误差信号}（error signal）。它可以计算一次、复用多次，这是反向传播算法效率的关键。


% Raster figure omitted from this study copy.



于是：
\begin{itemize}
    \item $\frac{\partial s}{\partial b} = \delta$（因为 $\frac{\partial z}{\partial b} = I$）
    \item $\frac{\partial s}{\partial W} = \delta^T x^T$（外积形式）
\end{itemize}

\textbf{计算 $\frac{\partial s}{\partial W}$}：


% Raster figure omitted from this study copy.



按照严格的 Jacobian 定义，$\frac{\partial s}{\partial W}$ 应该是一个 $1 \times nm$ 的行向量（因为 $s$ 是标量，$W$ 有 $nm$ 个参数）。但在实际工程中，我们采用\textbf{形状约定}（shape convention），将梯度整形为与参数\textbf{相同的形状}，以便直接做 SGD 更新 $W^{\text{new}} = W^{\text{old}} - \alpha \frac{\partial s}{\partial W}$。

最终结果：
$$\frac{\partial s}{\partial W} = \delta^T x^T \in \mathbb{R}^{n \times m}$$

这是 $\delta^T$（$n \times 1$）与 $x^T$（$1 \times m$）的\textbf{外积}。

\begin{warningbox}{Jacobian 形式 vs 形状约定}
数学上，Jacobian 矩阵有严格的定义和形状，链式法则在 Jacobian 形式下是矩阵乘法。但在工程实现中，为了方便 SGD 更新，我们通常将梯度``reshape''为与参数相同的形状。这两种约定混用可能导致困惑——Manning 教授建议：可以先用 Jacobian 做正确的数学推导，最后再 reshape 结果；或者始终按形状约定做，但要灵活使用转置。
\end{warningbox}


% Raster figure omitted from this study copy.



\subsection{外积梯度的直觉}

为什么 $\frac{\partial s}{\partial W} = \delta^T x^T$ 是对的？考虑单个权重 $W_{ij}$：它只参与计算 $z_i = \sum_k W_{ik} x_k + b_i$，即只连接输入 $x_j$ 到隐藏层 $z_i$。因此：

$$\frac{\partial s}{\partial W_{ij}} = \delta_i \cdot x_j$$

这恰好是外积 $\delta^T x^T$ 的第 $(i,j)$ 个元素。

\subsection{本章小结}

矩阵微积分的核心是：用 Jacobian 矩阵表示多输入多输出函数的导数，用矩阵乘法表示链式法则。对于神经网络的常见操作（线性层、激活函数、点积），Jacobian 都有简洁的封闭形式。关键概念是\textbf{上游梯度 $\delta$}——它在不同参数的梯度计算中被复用，避免重复计算。

%% ========== Part 6: 反向传播算法 ==========

\section{反向传播算法}

\subsection{反向传播的本质}

\begin{importantbox}{反向传播算法 = 两件事}
\begin{enumerate}
    \item \textbf{使用链式法则}：将复合函数的梯度分解为局部梯度的乘积
    \item \textbf{存储中间结果}：避免重复计算——同一个上游梯度被多个下游节点共享
\end{enumerate}
这就是反向传播算法的全部。
\end{importantbox}

\subsection{计算图}

为了系统化地实现反向传播，我们将计算过程表示为\textbf{计算图}（Computation Graph）：


% Raster figure omitted from this study copy.



计算图中：
\begin{itemize}
    \item \textbf{源节点}：输入变量（$x, W, b, u$）
    \item \textbf{内部节点}：运算操作（矩阵乘法、加法、激活函数、点积）
    \item \textbf{边}：传递操作的结果
\end{itemize}

\textbf{前向传播}（Forward Propagation）：沿计算图从左到右，依次计算每个节点的值。

\textbf{反向传播}（Backward Propagation / Backpropagation）：沿计算图从右到左，依次计算梯度。


% Raster figure omitted from this study copy.



\subsection{单节点的梯度传播规则}

对于计算图中的任意一个节点，设其计算为 $h = f(z)$：


% Raster figure omitted from this study copy.



$$\underbrace{\frac{\partial s}{\partial z}}_{\text{下游梯度}} = \underbrace{\frac{\partial s}{\partial h}}_{\text{上游梯度}} \times \underbrace{\frac{\partial h}{\partial z}}_{\text{局部梯度}}$$

这就是链式法则在计算图中的体现。当节点有多个输入时，对每个输入分别计算局部梯度，然后分别与上游梯度相乘。

\subsection{实例：$f(x,y,z) = (x+y) \cdot \max(y,z)$}

Manning 教授用一个具体的非神经网络的例子来演示反向传播的过程。


% Raster figure omitted from this study copy.



\textbf{前向传播}：
\begin{itemize}
    \item $a = x + y = 1 + 2 = 3$
    \item $b = \max(y, z) = \max(2, 0) = 2$
    \item $f = a \cdot b = 3 \times 2 = 6$
\end{itemize}

\textbf{局部梯度}：
\begin{itemize}
    \item 加法节点：$\frac{\partial a}{\partial x} = 1$, $\frac{\partial a}{\partial y} = 1$
    \item max 节点：$\frac{\partial b}{\partial y} = 1$（$y > z$），$\frac{\partial b}{\partial z} = 0$
    \item 乘法节点：$\frac{\partial f}{\partial a} = b = 2$，$\frac{\partial f}{\partial b} = a = 3$
\end{itemize}

\textbf{反向传播}（从右到左）：


% Raster figure omitted from this study copy.



\begin{enumerate}
    \item $\frac{\partial f}{\partial f} = 1$
    \item 乘法节点：$\frac{\partial f}{\partial a} = b = 2$，$\frac{\partial f}{\partial b} = a = 3$
    \item 加法节点：$\frac{\partial f}{\partial x} = \frac{\partial f}{\partial a} \cdot 1 = 2$，从加法到 $y$：$\frac{\partial f}{\partial a} \cdot 1 = 2$
    \item max 节点：到 $y$：$\frac{\partial f}{\partial b} \cdot 1 = 3$，到 $z$：$\frac{\partial f}{\partial b} \cdot 0 = 0$
    \item \textbf{$y$ 有两条路径}：$\frac{\partial f}{\partial y} = 2 + 3 = 5$
\end{enumerate}

最终结果：$\frac{\partial f}{\partial x} = 2$，$\frac{\partial f}{\partial y} = 5$，$\frac{\partial f}{\partial z} = 0$。

\textbf{验证}：将 $y$ 从 2 改为 2.1，则 $a = 3.1$，$b = 2.1$，$f = 3.1 \times 2.1 = 6.51$，变化量 $\approx 0.51 \approx 5 \times 0.1$，与梯度 5 吻合。

\subsection{分支处的梯度求和}


% Raster figure omitted from this study copy.



上例中 $y$ 同时参与了加法和 max 两个计算，因此它的总梯度是两条路径上梯度的\textbf{和}：

$$\frac{\partial f}{\partial y} = \frac{\partial f}{\partial a} \cdot \frac{\partial a}{\partial y} + \frac{\partial f}{\partial b} \cdot \frac{\partial b}{\partial y} = 2 \times 1 + 3 \times 1 = 5$$

这就是多元链式法则的体现。

\subsection{三种基本操作的梯度直觉}


% Raster figure omitted from this study copy.



\begin{importantbox}{三种基本操作的梯度行为}
\begin{itemize}
    \item \textbf{加法}（$+$）：\textbf{分发}梯度——上游梯度原封不动地传给两个输入
    \item \textbf{Max}：\textbf{路由}梯度——上游梯度全部传给较大的那个输入，另一个输入得到零
    \item \textbf{乘法}（$\times$）：\textbf{切换}梯度——上游梯度乘以\textbf{另一个}输入的值传给当前输入
\end{itemize}
\end{importantbox}

\begin{warningbox}{Max 操作导致梯度消失}
Max 操作对``输的那一方''梯度为零——这与 ReLU 的行为一致（$\text{ReLU}(x) = \max(0, x)$）。当 $x < 0$ 时，ReLU 的梯度为零，意味着该方向的信息完全被阻断。这从计算图的角度解释了 dying ReLU 现象。
\end{warningbox}

\subsection{高效反向传播：避免重复计算}


% Raster figure omitted from this study copy.



如果分别独立地计算 $\frac{\partial s}{\partial b}$、$\frac{\partial s}{\partial W}$、$\frac{\partial s}{\partial x}$、$\frac{\partial s}{\partial u}$，会发现它们共享大量的中间计算（上游梯度 $\delta$）。反向传播的核心效率在于：\textbf{按照拓扑排序的逆序一次性遍历计算图}，每个中间梯度只计算一次。

\begin{importantbox}{反向传播的计算复杂度}
如果正确实现，反向传播的时间复杂度与前向传播的时间复杂度\textbf{同阶}（Big-O 相同）。如果你的反向传播比前向传播慢很多，说明你在某处做了重复计算。
\end{importantbox}

\subsection{通用反向传播算法}


% Raster figure omitted from this study copy.



算法流程：

\textbf{前向传播}：
\begin{enumerate}
    \item 对计算图做\textbf{拓扑排序}，确保每个节点只依赖已计算的节点
    \item 按拓扑序遍历，对每个节点调用其 \texttt{forward} 方法计算输出值
\end{enumerate}

\textbf{反向传播}：
\begin{enumerate}
    \item 初始化输出梯度：$\frac{\partial s}{\partial s} = 1$
    \item 按\textbf{逆拓扑序}遍历节点
    \item 对每个节点：$\text{下游梯度} = \text{上游梯度} \times \text{局部梯度}$
    \item 对分支节点：将多条路径的梯度\textbf{求和}
\end{enumerate}

该算法适用于任意有向无环图（DAG），不要求网络是规整的层状结构。

\subsection{本章小结}

反向传播算法是链式法则在计算图上的高效实现。其核心是``一次前向、一次反向''，利用中间结果共享避免重复计算。对于任意 DAG 结构的计算图，反向传播的复杂度与前向传播同阶。

%% ========== Part 7: 自动微分与实现 ==========

\section{自动微分与框架实现}

\subsection{自动微分}

既然反向传播算法是如此系统化，能否让计算机\textbf{完全自动}完成？

\begin{knowledgebox}{自动微分的历史}
早期的深度学习框架 \textbf{Theano}（蒙特利尔大学开发）尝试了完全符号化的自动微分：给定前向计算的符号表达式，自动推导出反向传播的符号表达式。但这种方式过于重量级，难以灵活处理各种情况。
\end{knowledgebox}

现代深度学习框架（PyTorch、TensorFlow 等）采用了一种折中方案：

\begin{itemize}
    \item \textbf{框架负责}：管理计算图、执行拓扑排序、运行前向/反向遍历、传递上游梯度
    \item \textbf{用户负责}：为每种操作实现 \texttt{forward}（前向计算）和 \texttt{backward}（局部梯度计算）方法
\end{itemize}

\subsection{Forward/Backward API}


% Raster figure omitted from this study copy.



\begin{lstlisting}[caption={乘法门的 Forward/Backward 实现},language=Python]
class MultiplyGate(object):
    def forward(self, x, y):
        z = x * y
        self.x = x  # 缓存输入，供 backward 使用
        self.y = y
        return z

    def backward(self, dz):
        dx = dz * self.y  # 上游梯度 * 对方输入
        dy = dz * self.x
        return [dx, dy]
\end{lstlisting}

\begin{warningbox}{backward 依赖 forward 的缓存}
注意一个关键细节：\texttt{backward} 方法需要用到前向传播时的输入值（如乘法门中的 $x$ 和 $y$）。因此 \texttt{forward} 方法必须将这些值\textbf{缓存}起来（存为类的属性）。这就是为什么在 PyTorch 的 \texttt{autograd.Function} 中，\texttt{forward} 需要用 \texttt{ctx.save\_for\_backward()} 保存张量。
\end{warningbox}

实际上，PyTorch 已经预定义了大量常用操作的 forward 和 backward 实现。用户只需像搭积木一样组合这些操作，PyTorch 会自动管理计算图和梯度传播。这正是``高中生也能做深度学习项目''的原因。

\subsection{本章小结}

现代深度学习框架将反向传播的``基础设施''（计算图管理、梯度传播）自动化，用户只需为每个操作定义前向计算和局部梯度。常见操作（线性层、激活函数、损失函数等）都已预实现，使用者通常不需要手动编写 backward 方法。但理解底层原理对于调试和设计新操作至关重要。

%% ========== Part 8: 梯度检验 ==========

\section{梯度检验：数值验证}

\subsection{数值梯度}

如何验证你手动推导或实现的梯度是否正确？可以用\textbf{数值梯度}（Numerical Gradient）进行检验：

$$\frac{\partial f}{\partial x} \approx \frac{f(x + h) - f(x - h)}{2h}$$

其中 $h$ 是一个很小的数（通常取 $10^{-4}$）。


% Raster figure omitted from this study copy.



\begin{warningbox}{双侧差分 vs 单侧差分}
一定要用\textbf{双侧差分}（centered difference）$\frac{f(x+h) - f(x-h)}{2h}$，而不是单侧差分 $\frac{f(x+h) - f(x)}{h}$。双侧差分的误差是 $O(h^2)$，而单侧差分的误差是 $O(h)$，精度差了一个数量级。
\end{warningbox}

\subsection{为什么不直接用数值梯度训练？}

数值梯度计算简单，为什么不直接用它代替反向传播？因为它\textbf{极其缓慢}：需要对每一个参数单独做一次前向计算。如果模型有 $N$ 个参数，就需要 $2N$ 次前向传播（双侧差分）。而反向传播只需要\textbf{一次前向 + 一次反向}就能得到所有参数的梯度。

\begin{importantbox}{梯度检验的正确用法}
数值梯度仅用于\textbf{验证}你的 backward 实现是否正确，绝不用于训练。检验方法：
\begin{enumerate}
    \item 用反向传播计算梯度 $g_{\text{analytic}}$
    \item 用数值差分估计梯度 $g_{\text{numeric}}$
    \item 检查 $|g_{\text{analytic}} - g_{\text{numeric}}| < \epsilon$（通常 $\epsilon \sim 10^{-5}$）
\end{enumerate}
在 PyTorch 已预实现大部分操作的今天，梯度检验的需求减少了，但在实现\textbf{自定义操作}时仍然是必不可少的验证手段。
\end{importantbox}

\subsection{本章小结}

数值梯度提供了一种``暴力但可靠''的梯度估计方法，是验证解析梯度实现正确性的黄金标准。使用双侧差分可以获得更高的精度。但由于计算量与参数数量成正比，数值梯度只适合验证，不适合训练。

%% ========== 总结与延伸 ==========

\section{总结与延伸}

\subsection{讲者的核心总结}


% Raster figure omitted from this study copy.

\footnotetext{Slides 第85页（最后一页）。}

Manning 教授在课程结尾总结了以下要点：

\begin{enumerate}
    \item \textbf{反向传播 = 链式法则的高效实现}：前向传播计算函数值，反向传播计算梯度
    \item \textbf{理解底层原理很重要}：虽然现代框架（PyTorch 等）自动化了一切，但理解反向传播的数学有助于调试和设计新架构
    \item \textbf{梯度消失/爆炸}：后续讲解 RNN 时会看到，多层梯度连乘可能导致梯度指数级增长或衰减
    \item \textbf{实践中的平衡}：框架预定义了常用操作的 forward/backward，用户像拼积木一样组合即可
\end{enumerate}

\subsection{全课知识图谱}

\begin{center}
\begin{tikzpicture}[
    node distance=0.9cm,
    every node/.style={draw, rounded corners, minimum width=4cm, minimum height=0.7cm, font=\small, align=center},
    arrow/.style={-{Stealth[length=3mm]}, thick}
]
\node (nn) {神经网络基础};
\node (act) [below=of nn] {激活函数与非线性};
\node (sgd) [below=of act] {SGD 与梯度计算};
\node (mc) [below=of sgd] {矩阵微积分};
\node (bp) [below=of mc] {反向传播算法};
\node (impl) [below=of bp] {自动微分与框架实现};

\draw[arrow] (nn) -- (act);
\draw[arrow] (act) -- (sgd);
\draw[arrow] (sgd) -- (mc);
\draw[arrow] (mc) -- (bp);
\draw[arrow] (bp) -- (impl);

\node[draw=none, right=0.8cm of nn, text width=5.5cm, font=\scriptsize, align=left] {层、权重、偏置、矩阵表示};
\node[draw=none, right=0.8cm of act, text width=5.5cm, font=\scriptsize, align=left] {Sigmoid, tanh, ReLU, Swish, GELU};
\node[draw=none, right=0.8cm of sgd, text width=5.5cm, font=\scriptsize, align=left] {$\theta^{new} = \theta^{old} - \alpha\nabla J(\theta)$};
\node[draw=none, right=0.8cm of mc, text width=5.5cm, font=\scriptsize, align=left] {Jacobian、链式法则、$\delta$、Hadamard 积};
\node[draw=none, right=0.8cm of bp, text width=5.5cm, font=\scriptsize, align=left] {计算图、前向/反向传播、梯度求和};
\node[draw=none, right=0.8cm of impl, text width=5.5cm, font=\scriptsize, align=left] {forward/backward API、梯度检验};
\end{tikzpicture}
\end{center}

\subsection{关键 Takeaways}

\begin{importantbox}{五条核心原则}
\begin{enumerate}
    \item \textbf{非线性是必须的}：没有激活函数，再多层也只是一个线性变换
    \item \textbf{矩阵微积分 = 单变量微积分 + 矩阵}：Jacobian、链式法则在形式上与标量版完全类似
    \item \textbf{$\delta$（上游梯度）是核心概念}：它代表从损失函数回传到当前位置的``误差信号''，是多个参数梯度共享的部分
    \item \textbf{反向传播的效率来自复用}：按逆拓扑序一次遍历，每个中间梯度只计算一次
    \item \textbf{形状约定简化工程实现}：将梯度整形为与参数相同的形状，使 SGD 更新变为简单的减法
\end{enumerate}
\end{importantbox}

\subsection{拓展阅读}

\begin{itemize}
    \item CS224N 官方讲义（Lecture Notes）与矩阵微积分教程：\url{https://web.stanford.edu/class/cs224n/}
    \item Stanford Math 51 在线教材（线性代数与多变量微积分）：\url{http://web.stanford.edu/class/math51/textbook.html}
    \item Karpathy, ``Yes you should understand backprop''：\url{https://karpathy.medium.com/yes-you-should-understand-backprop-e2f06eab496b}
    \item Goodfellow et al., \textit{Deep Learning} Chapter 6: Deep Feedforward Networks
    \item Justin Johnson, ``Backpropagation for a Linear Layer'': \url{https://web.eecs.umich.edu/~justincj/teaching/eecs442/notes/linear-backprop.html}
    \item PyTorch 官方教程 -- Autograd: \url{https://pytorch.org/tutorials/beginner/blitz/autograd_tutorial.html}
\end{itemize}

\end{document}
